\medskip

On résout l'inéquation :
\begin{align*}
f(t) &\leqslant 44 \\
60\e^{-0,08t} + 23 &\leqslant 44 \\
60\e^{-0,08t} &\leqslant 21 \\
\e^{-0,08t} &\leqslant \dfrac{21}{60} \\
\e^{-0,08t} &\leqslant 0,35 \\
-0,08t &\leqslant \ln(0,35) \\
t &\geqslant \dfrac{\ln(0,35)}{-0,08}.
\end{align*}

Comme :
\[\dfrac{\ln(0,35)}{-0,08} \approx 13,12,\]
la température atteint donc $44\,^\circ\mathrm{C}$ environ $13,12$ minutes après la sortie de la machine.

Comme $f(13) \approx 44,2 > 44$ et $f(14) \approx 42,60 < 44$,
la première minute entière pour laquelle la température est inférieure ou égale à $44\,^\circ\mathrm{C}$ est la \textbf{14\up{e}} minute.
